\section{Detection by maps of fixed genus}\label{sec:genus}
Write $x_{h,m}$ for the scalar classes of
\cite[Definition~4.7]{Pardon}. The map $\Phi$ of \S\ref{sec:families}
will only have to distinguish scalar classes of multiplicity one.

\begin{theorem}\label{thm:independence}
The images $\Phi(x_{h,1})$, $h\ge0$, are linearly independent over $\QQ$.
\end{theorem}

We first reduce an arbitrary linear relation to a statement about a
single moduli space. Section~\ref{sec:degree} proves the required
nonvanishing there.

\subsection{The image cycle}
Let $Y\to B$ be a family as in \S\ref{sec:families}. Let $M_k$ be
the coarse moduli space of connected unmarked stable maps in the fibers,
of arithmetic genus $k$, with nonconstant image and vertical Chern number
zero. Write $z_k:M_k\to Z(Y/B)$ for the map taking the image cycle.
A component is counted with its degree over the normalization of its
image. Contracted components contribute zero.
\begin{lemma}\label{lem:proper}
The map $z_k$ is continuous and proper.
\end{lemma}
\begin{proof}
Gromov convergence gives convergence of the weighted currents and the
Hausdorff supports. Contracted components have image in the nonconstant part of the
image because the domain is connected. For properness, take a
sequence above a compact subset of $Z(Y/B)$ and pass to a subsequence whose
cycles converge. Work over one closed simplex, in a local target
trivialization about the limiting support. Hausdorff convergence puts
all the images in one relatively compact neighborhood. The masses of
the cycles bound the energies of the maps, with multiplicity included.
Indeed one can test against a cutoff Hermitian form of the limiting
parameter, uniformly positive on all nearby complex lines in this
neighborhood. For the fixed parametrizing surface in
Ivashkovich--Shevchishin's theorem take a closed oriented surface
of genus $k$. Each connected nodal domain of arithmetic genus $k$
is the quotient of such a surface by pinching disjoint circles,
as one sees by smoothing its nodes. The boundary condition in
their theorem is empty. The images stay in a common compact set,
and the structures on its neighborhood converge smoothly after
trivialization. Their theorem now gives a convergent subsequence
of stable maps. On smooth parts away from the nodes elliptic
regularity upgrades convergence to the smooth Gromov topology for
these smoothly converging structures. If the limit map were constant,
all nearby images would lie in a coordinate ball on which a fixed
exact two-form tames the nearby structures. Its integral over a
nonconstant closed curve is positive, whereas Stokes' theorem gives
zero. Thus the limit is nonconstant and its image cycle is the
prescribed limit. The empty cycle has a neighborhood disjoint from
all nonempty cycles, and the integer Chern number is locally constant
in the Gromov topology. These observations preserve the restrictions
defining $M_k$. The compactness argument applies in metrizable local
charts for the two topologies and over a finite union of simplices.
For compactness see \cite[Theorem~1.1]{IS}, and for the image-cycle map see \cite[Definition~2.14]{DIW}.
\end{proof}

\subsection{A regular map over the least genus}
Use a smooth curve $C$ of genus $h$ and a degree $h-1$ line $L$
outside the theta divisor. For $h=0$ take $\mathcal O(-1)$ and for
$h=1$ a nontrivial line of degree zero. Serre duality and
Riemann--Roch give $H^0(C,N)=H^1(C,N)=0$ for
$N=L\oplus K_C L^{-1}$. Every compact degree-one cycle in $N$ is a
section. Indeed it has one irreducible component mapping with degree
one to $C$, no compact curve lies in an affine fiber, and the
normalization gives the graph of a holomorphic section. Hence the
core is isolated in that degree and the point insertion represents
$x_{h,1}$, as in \S\ref{sec:local}.

For a stable map over this insertion, exactly one normalized
component maps with degree one to $C$ and it maps isomorphically.
Arithmetic genus equals the sum of the component genera plus the
first Betti number of the dual graph. It is at least $h$. In the
equality case the dual graph is a tree and every other vertex is a
contracted rational component. Any nontrivial such tree has a
contracted leaf, violating stability in the absence of markings.
Thus $z_h^{-1}([C])$ is the single embedded map, with trivial map
automorphisms in every genus. Its unparametrized deformation and
obstruction spaces are $H^0(C,N)$ and $H^1(C,N)$, since the
domain tangent complex cancels the tangential part of the target
complex. The map is regular with complex orientation $+1$.

In this splitting the whole degree-one cycle locus is the singleton
$[C]$. For both circles its unit represents the local class, by the
empty equivariant tuple at Borel stage zero. In particular
$a_{h,1}=x_{h,1}$. For $k<h$ the fiber of $z_k$ over this cycle is empty.

\begin{lemma}\label{lem:relation}
Suppose a finite relation among $\Phi(x_{h,1})$ has least genus $k$
with nonzero coefficient. There is a smooth family of almost-complex
six-manifolds over an open subset of $\CC^r$ for which the point
orientation at one regular map, in a product neighborhood of a
marked parameter, has zero image in $H_c^{2r}(M_k;\QQ)$.
\end{lemma}
\begin{proof}
Represent the sum on a finite disjoint
union of vertex targets, one per nonzero coefficient. In this
comparison diagram the filtered-quotient argument of Pardon's Lemma
3.2 applies. Take disjoint unions for common targets. For parallel
inclusions, attach a cylinder on the source base between their two
images. Use the corresponding pulled-back larger target near each
end and the source open target across the middle. Stacking three
cylinders permits simplicial attachment even when the endpoint
images overlap. The resulting target is Hausdorff since near the middle slice only the common open target is present, with
larger pieces on opposite sides. It is simplexwise smooth.
The two inclusions induce equal maps after attachment by the
constant-cylinder prism on compactly supported cochains and open
extension. A subdivision likewise has a triangulated product
concordance. Thus every finite vanishing relation is witnessed in
one family $X/B$ after these attachments. Distinct marked vertices
of the initial disjoint-union base remain distinct under its
inclusion~\cite[Lemmas 3.2--3.4]{Pardon}.

At each marked vertex $b$ attach a new closed interval to $B$ by
one endpoint only, and use the constant family with fiber $X_b$
on that interval. Its other endpoint $b'$ has an open neighborhood
in the enlarged base on which the family is a product.
The original base is a subcomplex, so its zero relation still
vanishes after inclusion. The compact point insertion crosses the
new edge by the interval relation for a constant family. Indeed
the insertion extended constantly over the edge has the difference
of its vertex insertions as its face boundary. Use distinct new
endpoints for the different marks. Relabel the enlarged complex
$B$ and use these endpoints as the marked parameters below.

Embed a subdivision of $B$ as a subcomplex in an oriented Euclidean
space and take a compact PL regular neighborhood, with $B$ in its
interior $W$. After subdivision there is a PL neighborhood retraction
onto $B$ that fixes $B$. Pull the family back along this retraction. The inclusion
of $B$ takes the zero relation to zero by the bivariant pushforward.
Equivalently, the dualizing complex of the neighborhood restricts
in the interior to $\mathbb Q[n]$, and the closed-base pushforward
has support on $B$ in that interior. One can use the pushforward
directly with target $W$, by properness of the closed inclusion,
so no restriction of compact supports across an arbitrary open set
is needed. Increase the ambient dimension if necessary. On $W$
the relation is therefore zero in $H_c^n(Z(X_W/W);\mathbb Q)$.
This is the base thickening used in the proof of Pardon's Theorem
5.15. It requires only a PL retraction for the present families~\cite[proof of Theorem 5.15]{Pardon}.

The closed-base pushforward for $i:B\hookrightarrow W$
uses the counit $i_*\omega_B=i_*i^!\omega_W\to\omega_W$.
Pull back this counit to the cycle space and use proper base change
for the closed inclusion of its restriction over $B$. This formula
also gives functoriality for a vertex followed by $i$, and commutes
with pulling back to the stable-map space by the proper map $z_k$.

Continuity of the retraction gives a Euclidean ball about each new
endpoint mapping into its constant-family neighborhood. Over this
ball the targets, cycle spaces and stable-map spaces are products.
The singleton cycle is open in the original open target and remains
open in the larger fiber by open embedding of the cycle spaces.
It is compact there as well. The oriented point Gysin map in the
ball is the external product of this singleton insertion and the
orientation at the marked parameter. Pullback by $z_k$ computes
the two factors separately. At a higher-genus mark the inverse image
of the singleton is empty by the genus bound above. At the least-genus mark it consists
of the one regular embedded map, so over the ball it is precisely
a copy of the ball. The pulled-back class is its point orientation
extended by zero, with a nonzero coefficient. This proves the reduction.

It remains to smooth the dependence of the targets on the parameter.
Let $P$ denote the compact triangulated PL neighborhood just constructed,
with interior $W$ open in Euclidean space. Increase its dimension to
an even integer $2r$. The family on $P$ is simplexwise smooth after
subdivision. Using the attached constant intervals, choose
disjoint Euclidean balls $B_{4\delta}(b)\Subset W$ around the
marked parameters on each of which the entire family is a product.
Use a common $\delta>0$ for these finitely many marks. The point
orientation on the distinguished stable-map graph maps to zero in
$H_c^{2r}(M_k)$ for a putative dependence. Subdivide to mesh
smaller than $\delta/4$.

For every vertex $v$ of the subdivision take its open star in $P$.
Choose a partition of unity $\lambda_v$ with closed supports contained
in these stars, smooth on the interior in the Euclidean smooth
structure. It can be obtained by restricting a smooth partition of
unity from an open Euclidean cover near $P$. Define $F:P\to P$ by
the barycentric combination of the vertices with weights
$\lambda_v(y)$. At $y$, the vertices whose closed supports contain
$y$ belong to the least simplex containing $y$. Since the cover is
finite, all other weights vanish in a whole neighborhood of $y$.
Consequently $F$ maps that neighborhood into a \emph{single} simplex
and its barycentric coordinates there are smooth in $W$. The
straight-line homotopy from $y$ to $F(y)$ stays in its simplex and
moves points by at most the mesh. In particular it preserves the
boundary of $P$.

Modify $F$ to fix the marked balls. Choose a smooth function
$\chi$ equal to zero on the balls of radius $\delta$ and one
outside the balls of radius $2\delta$. Replace $F(y)$ by
$(1-\chi(y))y+\chi(y)F(y)$, choosing $0\leq\chi\leq1$. This
map remains in the least simplex containing $y$, moves points by
less than $\delta/4$, and agrees with the old map near the boundary
of $P$. Where it changes, its source and image lie in one of the
balls of radius $3\delta$ with constant target family. Outside
the modifying balls the earlier factorization through one simplex
applies. Write $F$ for the modified map.

Set $W'=F^{-1}(W)$. The restriction $W'\to W$ is proper, as a base
change of the continuous map of compact spaces $P\to P$, and $W'$
lies in the Euclidean interior of $P$. Pull the target family back
along this map. It is a smooth submersion with smooth vertical
almost-complex structure on $W'$. On the modifying balls it is a
product. Elsewhere it is the pullback of a smooth simplex model.
On overlaps of simplex models the maps factor through the common
face and the prescribed identifications are smooth. The product
trivializations are the specified constant families on the attached
intervals, so they agree with the simplex models there. For the stable-map
spaces this is a Cartesian pullback, including their Gromov topologies,
so their induced map is also proper.

For a marked $b$, the preimage of $B_{\delta/2}(b)$ equals that
ball and $F$ is its identity. Indeed any preimage lies within
$3\delta/4$ of $b$ by the displacement bound. Consequently the
proper pullback of the point orientation on the distinguished graph
is the very same point orientation extended by zero. It still
vanishes in the global compactly supported group.
\end{proof}
